%% =============================================================================
%% sections/00_abstract_zh.tex — Chinese Abstract for Demo Paper
%% =============================================================================
纵向临床监护系统需要将异步高频生理波形与稀疏检验面板、结构化电子病历联合建模。现有多种多模态神经网络在回顾性队列上可获得较高区分度，但在传感器缺失、分布偏移和跨院区迁移下容易出现概率校准失真、静默退化与误报偏高。我们提出 \textbf{AuroraCard}，一套以校准为先的轻量多模态融合架构，围绕门控连续时间状态空间编码器与共形温度路由（conformal temperature routing）构建。在涵盖 184{,}200 例住院与门诊就诊的四中心评估基准上，AuroraCard 将 Macro-AUROC 提升至 \hlblue{96.4\%}（相比强 Transformer 基线提高 3.8 个百分点），预期校准误差（ECE）降低 \hlpink{61.5\%} 至 \hlpink{0.015}，单张工作站 GPU 上单个病人窗口推理时延为 \hlblue{14.2\,ms}。我们进一步给出结构化模态缺失条件下的有限样本覆盖保证，并公开模型权重、评测套件与可复现 LaTeX 报告模板。
